\figura[0.3]{fig1.png}

\begin{apartat}{1}
Un camp elèctric de més de $3,00\times 10^{6}\ \mathrm{V\,m^{-1}}$ provoca la ruptura dielèctrica de l'aire (els electrons són arrencats dels àtoms i en recombinar-se emeten llum). La descàrrega a través de l'aire causada per la ruptura dielèctrica s'anomena \textit{descàrrega en arc}. Un exemple familiar de descàrrega en arc és la descàrrega elèctrica que rebem quan toquem el pom metàl·lic d'una porta després d'haver caminat per una catifa en un dia sec. Calculeu, en aquest cas, la mínima diferència de potencial entre la mà i el pom de la porta si en el moment de la descàrrega elèctrica estan separats per $1,00\ \mathrm{mm}$.
\end{apartat}

\begin{apartat}{1}
Calculeu el treball que s'ha de fer perquè tres electrons que inicialment estaven molt separats quedin a $0,1\ \mathrm{nm}$ l'un de l'altre i configurin un triangle equilàter.
\end{apartat}

% DUBTE: l'original dona q_e = -1,60×10^{-9} C (errata per -1,60×10^{-19} C, el valor que fa servir la pauta).
\dades{%
  Càrrega de l'electró, $q_\mathrm{e} = -1,60\times 10^{-9}\ \mathrm{C}$.\\[2pt]
  $k = \dfrac{1}{4\pi\varepsilon_0} = 8,99\times 10^{9}\ \mathrm{N\,m^2\,C^{-2}}$.}
